\documentclass{exam}
\usepackage{amsmath, amssymb}

\pagestyle{headandfoot}
\firstpageheadrule
\runningheadrule
\firstpageheader{Prof. Tahvildar-Zadeh \\ Linear Algebra}{Homework 7}{Aadhithya Saravanan}
\runningheader{Linear Algebra \\ Homework 7}{}{Saravanan}
\firstpagefooter{}{}{}
\runningfooter{}{\thepage}{}

\printanswers

\begin{document}

\underline{Exercise 4.2.7}
\newline
\newline
Taking the determinant by expanding along the second row, we can ignore the second determinant, since $A_{2,2}$ is a zero entry. So,
\[\det(A)=(-1)^{2+1}\cdot(-1)\cdot\det(\tilde{A}_{2,1})+(-1)^{2+3}\cdot(-3)\cdot\det(\tilde{A}_{2,3})\]
\[=-(-1)\cdot(1\cdot0-2\cdot3)-(-3)\cdot(0\cdot3-1\cdot2)\]
\[=-6-6\]
\[=-12\]
\newline

\underline{Exercise 4.2.8}
\newline
\newline
Taking the determinant by expanding along the third row, we can ignore the third determinant, since $A_{3,3}$ is a zero entry. So,
\[\det(A)=(-1)^{3+1}\cdot(-1)\cdot\det(\tilde{A}_{3,1})+(-1)^{3+2}\cdot(3)\cdot\det(\tilde{A}_{3,2})\]
\[=(-1)\cdot(0\cdot5-2\cdot1)-3\cdot(1\cdot5-2\cdot0)\]
\[=2-15\]
\[=-13\]
\newline

\underline{Exercise 4.2.14}
\newline
\newline
Taking the determinant by expanding along the third row, we can ignore the second and third determinants, since $A_{3,2}$ and $A_{3,3}$ are zero entries. So,
\[\det(A)=(-1)^{3+1}\cdot(7)\cdot\det(\tilde{A}_{3,1})\]
\[=7\cdot(3\cdot0-4\cdot6)\]
\[=-168\]
\newline

\underline{Exercise 4.2.18}
\newline
\newline
We can use row reduction to simplify this determinant calculation.
\[
\begin{bmatrix}
    1 & -2 & 3 \\
    -1 & 2 & -5 \\
    3 & -1 & 2
\end{bmatrix}
\to
\begin{bmatrix}
    1 & -2 & 3 \\
    0 & 0 & -2 \\
    0 & 5 & -7
\end{bmatrix}
\to
\begin{bmatrix}
    1 & -2 & 3 \\
    0 & 5 & -7 \\
    0 & 0 & -2
\end{bmatrix}
\]
Let the original matrix be $A$ and the final matrix be $M$. The first row operation was of type 3, so the value of the determinant does not change. The second row operation was of type 1, so $\det(M)=-\det(A)$. Now, we can expand along row 3 of the matrix, ignoring the first and second determinants, since $M_{3,1}$ and $M_{3,2}$ are zero entries.
\[\det(M)=(-1)^{3+3}\cdot(-2)\cdot\det(\tilde{M}_{3,3})\]
\[=-2\cdot(1\cdot5-(-2)\cdot0)\]
\[=-10\]
Then, $\det(A)=-\det(M)=10$.
\newline

\underline{Exercise 4.2.23}
\newline
\newline
We can prove that the determinant of an upper triangular matrix is equal to the product of its diagonal entries by induction. An upper triangular matrix $A$ has $A_{i,j}=0$ whenever $i>j$. For the base case, let $N=1$. Since the determinant of a matrix of dimension $1\times1$ is defined as just the single entry, which is the only diagonal entry, the base case holds. Now, we assume for all upper triangular matrices $B$ of dimension $N-1\times N-1$ that
\[\det(B)=\prod_{i=1}^{N-1}B_{i,i}\]
Let $A$ be an upper triangular matrix of dimension $N\times N$. Then, since $A$ is upper triangular, $A_{N,j}=0$ for $j<N$. So, for the last row, the only potentially nonzero entry is $A_{N,N}$. Expanding the determinant along the $N$th row, we have
\[\det(A)=(-1)^{N+N}\cdot(A_{N,N})\cdot\det(\tilde{A}_{N,N})\]
$N+N$ is twice a natural number, so it is even, and we have
\[\det(A)=(A_{N,N})\cdot\det(\tilde{A}_{N,N})\]
Now, $\tilde{A}_{N,N}$ is a matrix of dimension $N-1\times N-1$. It is also upper triangular, as it is just $A$ without the $N$th row and $N$th column. By the induction hypothesis, \[\det(\tilde{A}_{N,N})=\prod_{i=1}^{N-1}A_{i,i}\]
So, substituting this into the equation with $\det(A)$, we have
\[\det(A)=(A_{N,N})\cdot\prod_{i=1}^{N-1}A_{i,i}\]
\[=\prod_{i=1}^{N}A_{i,i}\]
This is equal to the product of the diagonal entries of a matrix $A$ of dimension $N\times N$. Therefore, the determinant of an upper triangular matrix is the product of its diagonal entries.
\newline

\underline{Exercise 4.3.12}
\newline
\newline
We are given that $QQ^t=I$. Taking the determinant of each side, we have $\det(QQ^t)=\det(I)$. By Theorem 4.7, $\det(QQ^t)=\det(Q)\det(Q^t)$. By Theorem 4.8, $\det(Q^t)=\det(Q)$. Substituting these back, we have $\det(Q)\det(Q)=\det(Q)^2=\det(I)$. Since $\det(I)=1$ by Example 4.2.4, $\det(Q)^2=1$. This implies $\det(Q)=\pm1$.
\newline

\underline{Exercise 4.3.15}
\newline
\newline
If $A$ and $B$ are similar, there exists an invertible matrix $Q$ such that $B=Q^{-1}AQ$. Taking the determinant of each side, we have $\det(B)=\det(Q^{-1}AQ)$. By Theorem 4.7, $\det(Q^{-1}AQ)=\det(Q^{-1})\det(A)\det(Q)$. Let $q=\det(Q)$. By the corollary to Theorem 4.7, $\det(Q^{-1})=\frac{1}{\det(Q)}=\frac{1}{q}$. Substituting these back, we have $\det(B)=\frac{1}{q}\det(A)q=\det(A)$. Therefore, if $A$ and $B$ are similar, $\det(A)=\det(B)$.
\newline

\underline{Exercise 4.3.24}
\newline
\newline
Let $A_n$ represent the given matrix in $M_{n\times n}(F)$. We are going to prove that \[\det(A_n+tI_n)=t^n+a_{n-1}t^{n-1}+\dots+a_1t+a_0\] using induction on $n$. Let $n=1$. Then, $A_n$ has only the entry $a_0$. $tI_n$ has only the entry $t$, so $A_n+tI_n$ has only the entry $t+a_0$. The definition of the determinant says the determinant of matrix of dimension $1\times1$ is just the entry, so $\det(A_1+tI_1)=t+a_0$. This matches the formula, as it says the determinant for $n=1$ is $t^1+a_0=t+a_0$. Let $n>1$ and assume the formula is true for $n-1$. That is, \[\det(A_{n-1}+tI_{n-1})=t^{n-1}+a_{n-1}t^{n-2}+\dots+a_2t+a_1\] This is because we are using $A_{n-1}$ as $A_n$ without the first row and first column. Now, we show that the equation is true for $n$. Let $M=A_n+tI_n$. Expanding along the first row of $M$, we have the only potentially nonzero terms as $t$ in the $(1,1)$ position and $a_0$ in the $(1,n)$ position. Then, \[\det(M)=(-1)^{1+1}\cdot t\cdot\det(\tilde{M}_{1,1})+(-1)^{1+n}\cdot a_0\cdot\det(\tilde{M}_{1,n})\] $\tilde{M}_{1,1}$ is the same type of matrix as $M$ with one less row and one less column, so we can apply the induction hypothesis to get \[\det(\tilde{M}_{1,1})=t^{n-1}+a_{n-1}t^{n-2}+\dots+a_2t+a_1\]$\tilde{M}_{1,n}$ has the first row and last column removed, so its only potentially nonzero entries are the $-1$s on the diagonal and the $t$s above the diagonal. So, it is upper triangular and its determinant is the product of the diagonal entries, which are all $-1$. There are $n-1$ terms, so $\det(\tilde{M}_{1,n})=(-1)^{n-1}$. Substituting this back, we have \[\det(M)=(-1)^{1+1}\cdot t\cdot(t^{n-1}+a_{n-1}t^{n-2}+\dots+a_2t+a_1)+(-1)^{1+n}\cdot a_0\cdot(-1)^{n-1}\]
\[=1\cdot(t^n+a_{n-1}t^{n-1}+\dots+a_2t^2+a_1t)+(-1)^{2n}a_0\]
\[=t^n+a_{n-1}t^{n-1}+\dots+a_1t+a_0\]
\newline

\underline{Exercise 4.4.2.c}
\newline
\newline
\[\det(A)=(2+i)(3-i)-(-1+3i)(1-2i)\]
\[=(7+i)-(5+5i)\]
\[=2-4i\]
\newline

\underline{Exercise 4.4.3.c}
\newline
\newline
Taking the determinant by expanding along the second column, we can ignore the second determinant, since $A_{2,2}$ is a zero entry. So,
\[\det(A)=(-1)^{1+2}\cdot(1)\cdot\det(\tilde{A}_{1,2})+(-1)^{3+2}\cdot(3)\cdot\det(\tilde{A}_{3,2})\]
\[=-1\cdot(-1\cdot0-(-3)\cdot2)-3\cdot(0\cdot(-3)-2\cdot(-1))\]
\[=-6-6\]
\[=-12\]
\newline

\underline{Exercise 4.4.3.e}
\newline
\newline
Taking the determinant by expanding along the third row, we can ignore the third determinant, since $A_{3,3}$ is a zero entry. So,
\[\det(A)=(-1)^{3+1}\cdot(3)\cdot\det(\tilde{A}_{3,1})+(-1)^{3+2}\cdot(4i)\cdot\det(\tilde{A}_{3,2})\]
\[=3\cdot((1+i)\cdot(1-i)-2\cdot0)-4i\cdot(0\cdot(1-i)-2\cdot(-2i))\]
\[=6+16\]
\[=22\]
\newline

\end{document}